\section*{Capítulo \ref{ch:g9:genes-and-dna} --- Los genes y el ADN}

\begin{solution}{exo:g9:genes-and-dna:1}
Un \omterm{def:g9:genes-and-dna:gene}{gen} es un trozo de \omterm{def:g9:chromosomes:chromosomes}{cromosoma} que lleva las instrucciones de un
\omterm{def:g9:heredity-and-traits:traits}{rasgo}, en una dirección fija; los \omterm{def:g9:genes-and-dna:gene}{alelos} son sus versiones --- las A, B
y O del \omterm{def:g9:genes-and-dna:gene}{gen} del grupo sanguíneo.
\end{solution}

\begin{solution}{exo:g9:genes-and-dna:2}
Porque los \omterm{def:g9:chromosomes:chromosomes}{cromosomas} vienen en parejas: una copia de cada \omterm{def:g9:genes-and-dna:gene}{gen} viaja en
cada miembro --- una de la madre y otra del padre, en la misma
dirección.
\end{solution}

\begin{solution}{exo:g9:genes-and-dna:3}
El orden de las letras químicas a lo largo de la hebra. Cuatro clases
de letra, y unos tres mil millones de ellas en una baraja humana.
\end{solution}

\begin{solution}{exo:g9:genes-and-dna:4}
Las hebras son complementarias letra a letra: separadas, cada una
reconstruye a su compañera, y una escalera se convierte en dos
escaleras idénticas --- una copia completa por \omterm{def:g5:first-look-at-cells:cell}{célula} hija.
\end{solution}

\begin{solution}{exo:g9:genes-and-dna:5}
Un error de copia poco frecuente --- letras cambiadas que crean un
\omterm{def:g9:genes-and-dna:gene}{alelo} nuevo. Casi siempre neutra, a veces dañina y de vez en cuando
útil.
\end{solution}

\begin{solution}{exo:g9:genes-and-dna:6}
Todos los \omterm{def:g6:exploring-our-environment:components}{seres vivos} escriben y leen sus \omterm{def:g9:genes-and-dna:gene}{genes} con las mismas letras
del \omterm{def:g9:genes-and-dna:dna}{ADN} y el mismo código. Consecuencias: la prueba más fuerte de un
origen común --- y que mudar \omterm{def:g9:genes-and-dna:gene}{genes} entre \omterm{def:g5:biodiversity:species}{especies} funcione, porque el
idioma coincide en todas partes.
\end{solution}

\begin{solution}{exo:g9:genes-and-dna:7}
Cada progenitor lleva marrón con azul --- el marrón expresado y el azul
callado. Cada uno reparte azul con una posibilidad entre dos; los
repartos son independientes, así que azul con azul cae con una
posibilidad entre cuatro --- el hijo de \omterm{def:g1:the-five-senses:sense}{ojos} azules.
\end{solution}

\begin{solution}{exo:g9:genes-and-dna:8}
Grupo A: A con A, o A con O. Grupo B: B con B, o B con O. Grupo AB: A
con B --- los dos expresados. Grupo O: O con O.
\end{solution}

\begin{solution}{exo:g9:genes-and-dna:9}
Las dos tienen la biblioteca entera --- los mismos 46 hilos. La \omterm{def:g8:nervous-communication:neuron}{neurona}
lee los \omterm{def:g9:genes-and-dna:gene}{genes} de la maquinaria nerviosa y deja cerrados los capítulos
de la queratina de la \omterm{def:g1:the-five-senses:sense}{piel}; la \omterm{def:g5:first-look-at-cells:cell}{célula} de la \omterm{def:g1:the-five-senses:sense}{piel} lee lo contrario. La
especialización es una política de préstamo, no otra biblioteca.
\end{solution}

\begin{solution}{exo:g9:genes-and-dna:10}
Machaca la fruta (\omterm{def:g5:first-look-at-cells:cell}{células} abiertas); agua con sal y lavavajillas (las
\omterm{def:g6:cell-unit-of-life:parts}{membranas} --- las envolturas de la \omterm{def:g5:first-look-at-cells:cell}{célula} y del \omterm{def:g6:cell-unit-of-life:parts}{núcleo} --- disueltas,
el \omterm{def:g9:genes-and-dna:dna}{ADN} liberado y mantenido en disolución); cuela (fuera los restos);
alcohol frío echado por encima (el \omterm{def:g9:genes-and-dna:dna}{ADN} se sale de la disolución en la
frontera): hilos blanquecinos que se enrollan --- la molécula misma.
\end{solution}

\begin{solution}{exo:g9:genes-and-dna:11}
Nivel de familia: la barbilla reaparece siguiendo las líneas de la
reproducción, llevada sin mostrarse durante una generación. Nivel de
\omterm{def:g9:chromosomes:chromosomes}{cromosoma}: su determinante viaja en un hilo, reducido a la mitad en los
\omterm{def:g8:sexual-reproduction:gametes}{gametos} y restaurado en la \omterm{def:g5:flower-to-fruit:steps}{fecundación}, repartido a unos hijos y a
otros no. Nivel de molécula: es un \omterm{def:g9:genes-and-dna:gene}{alelo} --- una ortografía de algún
\omterm{def:g9:genes-and-dna:gene}{gen} que da forma a la cara --- copiado con fidelidad durante el siglo
que ha pasado desde la \omterm{prop:g9:genes-and-dna:explains}{mutación} que lo escribió por primera vez.
\end{solution}

\begin{solution}{exo:g9:genes-and-dna:12}
La \omterm{prop:g9:genes-and-dna:explains}{mutación} escribe: los errores de copia crean ortografías nuevas ---
nivel de molécula, este capítulo. El barajado reparte: la reducción a
la mitad y la \omterm{def:g5:flower-to-fruit:steps}{fecundación} le entregan a cada hijo una combinación nueva
--- nivel de \omterm{def:g9:chromosomes:chromosomes}{cromosoma}, \cref{ch:g9:chromosomes} (y
\cref{prop:g8:sexual-reproduction:shuffle}). La expresión enseña: de
las dos copias repartidas, qué se muestra y qué viaja callado --- la
regla de las dos copias de este capítulo. Solo la \omterm{prop:g9:genes-and-dna:explains}{mutación} crea
ortografías de verdad nuevas; las otras dos recombinan y revelan lo que
la \omterm{prop:g9:genes-and-dna:explains}{mutación} escribió una vez.
\end{solution}

\begin{solution}{pb:g9:genes-and-dna:1}
\textbf{1.} Una \omterm{prop:g9:genes-and-dna:explains}{mutación}: un error de copia en el \omterm{def:g9:genes-and-dna:gene}{gen} del pigmento,
una vez, en la línea de \omterm{def:g8:sexual-reproduction:gametes}{gametos} de algún antepasado --- letras
cambiadas, instrucciones estropeadas, y la falta de ortografía copiada
con fidelidad desde entonces.

\textbf{2.} Porque las instrucciones que funcionan de P construyen el
pigmento sin importar el compañero callado: p viaja \emph{llevado sin
mostrarse} --- la palabra del álbum, ahora molecular.

\textbf{3.} Con p en los dos miembros, no hay instrucciones que
funcionen en ninguna de las dos direcciones: la maquinaria del pigmento
--- el producto del \omterm{def:g9:genes-and-dna:gene}{gen} --- sencillamente no se construye nunca. El
\omterm{def:g9:heredity-and-traits:traits}{rasgo} es la ausencia de una ortografía que funcione.

\textbf{4.} Que sus \omterm{def:g9:genes-and-dna:gene}{genes} del pigmento son reconociblemente el mismo
\omterm{def:g9:genes-and-dna:gene}{gen} --- secuencias parecidas en el mismo código --- copias parientes de
un solo texto antepasado, estropeadas de manera independiente en la
historia de cada \omterm{def:g5:biodiversity:species}{especie}.

\textbf{5.} Los cuatro repartos: P con P, P con p, p con P y p con p
--- igual de probables. La obra sin pigmento: una posibilidad entre
cuatro.

\textbf{6.} El salto de generación: dos portadores que enseñaban
pigmento transmitieron lo que no mostraban, y las ortografías
escondidas se encontraron --- el \cref{ex:g9:heredity-and-traits:skip}
con su mecanismo entero.

\textbf{7.} Todos llevan una p --- la aportación entera de su madre en
esa dirección es p. Ninguno enseña la condición: la P de la pareja
construye el pigmento en todos ellos.

\textbf{8.} La p de la abuela viajaba en un \omterm{def:g9:chromosomes:chromosomes}{cromosoma} de una pareja; la
reducción a la mitad de su \omterm{def:g8:sexual-reproduction:gametes}{gameto} repartió ese miembro al \omterm{def:g4:animal-reproduction:sexual}{óvulo} o al
\omterm{def:g8:reproductive-systems:male}{espermatozoide} que hizo al progenitor portador; la reducción a la mitad
del progenitor lo repartió otra vez al \omterm{def:g8:sexual-reproduction:gametes}{gameto} que hizo al nieto ---
hilo, \omterm{def:g8:sexual-reproduction:gametes}{gameto} y \omterm{prop:g8:sexual-reproduction:core}{cigoto}, dos veces seguidas.

\textbf{9.} Familia: dos padres de aspecto corriente, un hijo de una
entre cuatro, y el hecho de llevarlo demostrado por el \omterm{def:g2:born-grow-die:cycle}{nacimiento}.
\omterm{def:g9:chromosomes:chromosomes}{Cromosoma}: los dos miembros de una pareja llevando P y p, reducidos a
la mitad y recombinados. Molécula: un \omterm{def:g9:genes-and-dna:gene}{gen}, un \omterm{def:g9:genes-and-dna:gene}{alelo} que funciona y otro
mal escrito, y la diferencia de ortografía decidiéndolo todo.

\textbf{10.} De todos los juicios: en un portador, lo estropeado de p
es invisible --- la P de la pareja lo tapa, así que nada de la obra del
portador ``informa'' de la falta de ortografía. Lo que no se puede ver
no se puede reparar ni quitar de en medio; las ortografías escondidas
duran indefinidamente.

\textbf{11.} Su \omterm{def:g9:heredity-and-traits:traits}{rasgo} (sin pigmento) es heredado --- p con p --- pero
su \emph{consecuencia} se vive en un ambiente lleno de sol: el pigmento
que falta era la crema solar del cuerpo, así que el ambiente escribe
sobre ella más que sobre sus padres protegidos --- la \omterm{def:g9:heredity-and-traits:heredity}{herencia} pone el
margen y el mundo lo rellena.

\textbf{12.} Lo que copia: la doble hebra, con cada mitad
reconstruyendo a su compañera en cada división celular y en cada
generación. Lo que se copió mal una vez: unas pocas letras del \omterm{def:g9:genes-and-dna:gene}{gen} del
pigmento, en un \omterm{def:g8:sexual-reproduction:gametes}{gameto} antiguo. Por qué dura el error: la fotocopiadora
es fiel --- copia lo que hay, y las faltas de ortografía con la misma
fidelidad que el sentido.
\end{solution}
